% Capstone 07 – Pollard's rho for the ECDLP. Port of «Мини-курсовая - ро-алгоритм Полларда
% решения ECDLP». The numbers of «Требования к реализации» come from requirements.tex,
% generated by `make bench CAPSTONE=pollard-rho-ecdlp` (see bench/README.md).
\documentclass{ca-capstone}
\capstonerefs{sagemath,hps}
% Generated by `make bench CAPSTONE=pollard-rho-ecdlp` (taskgen.capstone report) from
% bench/results.json: reference time x slack on the reference machine.
% Do not edit: change bench/config.yaml and re-run the benchmark.
\newcommand{\ReqSlack}{10}
\newcommand{\ReqPBitsLo}{50}
\newcommand{\ReqPBitsHi}{100}
\newcommand{\ReqOrderBits}{40}
\newcommand{\ReqSeconds}{200}

\begin{document}
\capstonetitle{Решение задачи ECDLP \texorpdfstring{$\boldsymbol{\rho}$}{ρ}-алгоритмом Полларда}

\goal
Требуется реализовать $\rho$-алгоритм Полларда для решения задачи дискретного
логарифмирования в группе точек на эллиптической кривой $E(\GF{p})$.

\functionality
\begin{enumerate}
  \item Возможность задавать значения параметров поиска: координаты стартовой точки и
    максимального числа итераций цикла поиска коллизии.
  \item Вычисление порядка основания дискретного логарифма,
    $\ord(P)$\footnote{для чего нам необходима эта величина?}.
  \item Нахождение решения \emph{ECDLP} и проверка корректности решения.
\end{enumerate}

\requirements
Требования к реализации:
\begin{enumerate}
  \item корректность решения задачи \emph{ECDLP};
  \item приемлемая скорость работы (см.\ комментарий ниже).
\end{enumerate}

\textbf{Корректность работы} можно проверять, самостоятельно генерируя задачи, либо же
прибегая к помощи сторонних решателей, например, функций библиотеки~\cite{sagemath}.

В данной работе важна не скорость, а корректность реализации и соответствие требованиям.
Однако тестовые примеры с числами размеров примерно
\textbf{\ReqPBitsLo--\ReqPBitsHi~бит} (при порядке основания $\ord(P)$ до
\textbf{\ReqOrderBits~бит}) должны выполняться за приемлемое время "--- не более
\ReqSeconds~секунд.

\details
Так как цель работы "--- решение \emph{ECDLP}, в своей реализации \textbf{можно пользоваться
библиотеками} для работы с точками на эллиптической кривой, а также пакетами линейной
алгебры.

Рекомендуется использовать язык \emph{Python} и систему компьютерной алгебры
\emph{SageMath}~\cite{sagemath}.

Желательно сохранять промежуточные результаты решения \emph{ECDLP} $\rho$-алгоритмом
Полларда, такие как значения коэффициентов $\alpha_i, \beta_i, \gamma_i, \delta_i$,
сформированное сравнение соответствующих найденной коллизии степеней, все возможные решения
этого сравнения и~проч.

Платформо-независимый, \textbf{переносимый код} писать \textbf{не требуется}, но во время
сдачи должна быть возможность проверить работу.

\schedule
Работа выполняется \textbf{индивидуально}. Срок выполнения "--- \textbf{месяц}.

Работа должна быть ваша: реализация должна быть \textbf{разработана и написана вами}, а не
собрана из кусков чужих работ.

Желательно хранить код в \emph{git}-репозитории на \emph{github}, \emph{bitbucket} или
подобных платформах. \textbf{Документация не требуется}.

\remarks
\begin{enumerate}
  \item Порядок точки $P$ можно вычислять, применяя функции библиотеки~\cite{sagemath} (или
    аналогичной), однако собственная реализация некоторого алгоритма поиска порядка точки
    (например, метода Шенкса"--~Местре) будет весомым плюсом вашей работы.
\end{enumerate}

\references
\end{document}
